% ---
% id: EX-CH01-006
% titre: Nombres amicaux
% chapitre_id: 1
% chapitre_nom: Nombres entiers & divisibilité
% annee_scolaire: 9H
% difficulte: 3
% tags: [divisibilite, diviseurs]
% auteur: Bussy D.
% ---

\begin{exercice}{EX-CH01-006}{Nombres amicaux}

\begin{mdframed}[backgroundcolor=vlmgray, innerleftmargin=10pt, innerrightmargin=10pt, innertopmargin=6pt, innerbottommargin=6pt, skipabove=-5pt, skipbelow=5pt]
Deux nombres sont dits \textbf{amicaux} si chacun est la somme des diviseurs stricts de l'autre.
\end{mdframed}

\bigskip
Vérifie les paires de nombres amicaux suivantes :

\medskip
\textbf{A.}

\medskip
La paire $(220 \;;\; 284)$ fut découverte par les Pythagoriciens. Elle fut considérée comme symbole de l'amitié.

\vspace{20pt}
\textbf{B.}

\medskip
La paire $(17296 \;;\; 18416)$ fut découverte par Fermat en 1636.

\end{exercice}

\begin{solution}{EX-CH01-006}

\textbf{A.} Paire $(220 \;;\; 284)$

\medskip
\begin{tabular}{ll}
$D^*(220)$ & $= \{1, 2, 4, 5, 10, 11, 20, 22, 44, 55, 110\}$ \\[4pt]
           & $1+2+4+5+10+11+20+22+44+55+110 = 284$ \quad $\checkmark$ \\[8pt]
$D^*(284)$ & $= \{1, 2, 4, 71, 142\}$ \\[4pt]
           & $1+2+4+71+142 = 220$ \quad $\checkmark$
\end{tabular}

\medskip
$\Rightarrow$ $220$ et $284$ sont bien amicaux.

\bigskip

\textbf{B.} Paire $(17296 \;;\; 18416)$

\medskip
\begin{tabular}{ll}
$D^*(17296)$ & $= \{1, 2, 4, 8, 16, 23, 46, 47, 92, 94, 184, 188, 376, 752, 1081, 2162, 4324, 8648\}$ \\[4pt]
             & $\sum D^*(17296) = 18416$ \quad $\checkmark$ \\[8pt]
$D^*(18416)$ & $= \{1, 2, 4, 8, 16, 1151, 2302, 4604, 9208\}$ \\[4pt]
             & $\sum D^*(18416) = 17296$ \quad $\checkmark$
\end{tabular}

\medskip
$\Rightarrow$ $17296$ et $18416$ sont bien amicaux.

\end{solution}